\subsubsection{Incremental value of larger UPS batteries}
The UPS sensitivity addresses whether additional energy capacity could be justified partly by its flexibility value. Increasing the central 600 kWh capacity to 750 and 900 kWh produces additional annual electricity-cost savings of GBP~1,292.28 and GBP~2,481.41, respectively. With the reserve fraction fixed at 50\%, the additions of 150 and 300 nameplate kWh increase the operational energy window by 75 and 150 kWh. Power ratings remain unchanged.

For a constant annual incremental saving $\Delta S$, additional annual non-electricity cost $O$, real discount rate $r$ and life $n$, the break-even incremental investment is
\begin{equation}
I_{\mathrm{BE}}=\frac{\Delta S-O}{\mathrm{CRF}(r,n)},\qquad
\mathrm{CRF}(r,n)=\frac{r(1+r)^n}{(1+r)^n-1}.
\label{eq:ups-investment}
\end{equation}
The capital recovery factor annualises initial expenditure~\cite{sam_crf}. At an illustrative 8\% real rate over ten years, with $O=0$, the additional savings support investments of approximately GBP~8,671 and GBP~16,650, or GBP~57.81 and GBP~55.50 per added nameplate kWh. These are expenditure thresholds, not installed battery prices or an unconditional case for enlargement. Installation, degradation, O\&M and replacement costs would need to be covered, and the retained single-year savings are not demonstrated lifetime cash flows: initial and terminal energy differ between capacities. Supplementary Section~S6 gives the assumptions and alternative rate/life combinations. No revenue for the physical flexibility events is added to these electricity-cost savings.
